% LSPV-4::ADAPTIVEMOMENT1::2026-10-15_0840KST
% Full eight-page manuscript and reproducibility bundle preserved as local artifacts.
% Frozen ideas SHA256 c66297126ddda8c67a44bbad3acccc4eb4753f0d11c5383d66bafb2f0bf0d46a
\documentclass[11pt]{article}
\usepackage{amsmath,amssymb,amsthm,hyperref}
\newtheorem{theorem}{Theorem}
\newtheorem{lemma}[theorem]{Lemma}
\title{ADAPTIVEMOMENT1: Exact Packed Low-Degree Extension for Two Adaptive Indexed Reads}
\author{LSPV-4, Nonrelativizing Simulation}
\date{15 October 2026, 08:40 Asia/Seoul}
\begin{document}\maketitle
\noindent\textbf{STATUS.} \textbf{PROVED/EXACT} for the restricted explicit source grammar described below, independent review pending. Eight final computational batches: 10,284 assertions, zero mismatches, status \textbf{COMPUTATIONAL}. No arbitrary decider-to-source compiler, no complete Williams--Cook--Mertz graph time bound, no unconditional nonrelativizing SAT separation or nonalgebrization theorem. $P\ne NP$ remains \textbf{OPEN}.

\section{IDEAS PROPOSED BEFORE LITERATURE}
Original proposals were frozen at SHA-256 \texttt{c66297126ddda8c67a44bbad3acccc4eb4753f0d11c5383d66bafb2f0bf0d46a} before new external literature search, separately from revisions.
\par\noindent\textbf{CANDIDATE A.} Compute the canonical packed extension of $X_{B_{X_A}}$ using exact repeated-data-bit collision moments. Missing premise: cheap within-pack two-bit moments. Falsifier: off-grid disagreement with full interpolation.
\par\noindent\textbf{CANDIDATE B.} Trace-basis products create one top-degree reduction term, yielding rank-one within-pack off-diagonal corrections. Missing premise: no other term of degree at least $q$. Falsifier: direct two-bit tensor-interpolation disagreement.
\par\noindent\textbf{CANDIDATE C.} Compute all address-fiber weights without an $m$-entry table and stream data packs twice. Missing premise: compatible $q$, fields, and simultaneous bit cost. Falsifier: $m^2$ scans or retained $m$-sized state.
\par\noindent\textbf{CANDIDATE D.} Contrast one explicit local source with an unrestricted Boolean oracle gate and completion-time fresh samples. Missing premise for separation: a generic source-level compiler and bounded whole-graph call schedule. Falsifier: an oracle off-grid shortcut with no oracle-extension advice or a compulsory offline SAT preprediction obligation.

\section{LITERATURE CHECK AND REVISIONS}
Williams, \emph{Simulating Time With Square-Root Space} (STOC 2025), Theorems 1.1 and 3.1, provides a multitape space simulation, not a near-original-time simulation. Remark 1.6 explicitly excludes the general random-access machine model. Section 5 identifies naive local low-degree extension evaluation as a bottleneck. Theorem 5.1 preserves an oracle-query-length-restricted variant while the full theorem is nonrelativizing \cite{Wil}. Cook--Mertz tree evaluation controls space, not an adequate fast-time schedule \cite{CM}. Henzinger--Pyne--Ragavan's 2026 catalytic tree evaluation uses a different memory model \cite{HPR}. The trace-dual and linearized-polynomial ingredients are standard finite-field algebra \cite{LN}, not a general SAT result. Revisions: A and B survive after treating both same-index and same-pack correlations; C survives with $g=k+5,q=32m$ and two charged data passes. D is retained as a scoped oracle diagnostic and rejected as a complete separation.

\section{CONSTRUCTION OR PROOF ATTEMPT}
Let $m=2^k\ge4$, $g=k+5$, $q=2^g=32m$, $K=\mathbb F_q$ in a polynomial basis, and $E=\mathbb F_{q^2}$. There are three full $g$-bit address words $A,B_0,B_1$ and an $m$-bit data array $X$, padded to whole $g$-bit packs. An address denotes the integer encoded by its lowest $k$ bits. Set
\[
 f_m(A,B_0,B_1,X)=X_{B_{X_A}},\quad
 a=3g+g\lceil m/g\rceil\le m+4g-1,\quad p=3+\lceil m/g\rceil.
\]
Since $2^k\ge4(k+5)/7$ for $k\ge2$, $q\ge4(a+1)$, and $p(q-1)<q^2-1$. This matches Williams's strong packed-degree regime. The Boolean source is computable by a fixed multitape program in $O(m)$ time, including physical sequential index access. For $\alpha\in K$, the packed-grid Lagrange polynomial is $\delta_\alpha(Z)=1+(Z+\alpha)^{q-1}$, with $\delta_\alpha(\beta)=[\alpha=\beta]$ on $K$. The canonical extension $\widetilde f_m(z)$ is the unique polynomial of degree less than $q$ in each of its $p$ coordinates agreeing with $f_m$ on $K^p$.

Let $\theta_0,\dots,\theta_{g-1}$ be the trace-dual basis, so coordinate bit $i$ on $K$ is $\mathrm{Tr}_{K/\mathbb F_2}(\theta_i\alpha)$. Define
\[
 L_i(Z)=\sum_{r=0}^{g-1}(\theta_i Z)^{2^r},\quad c_i=\theta_i^{2^{g-1}},\quad \eta(Z)=Z^q+Z.
\]
\begin{lemma}[Canonical two-bit moment]\label{lem:pair}
For two bit positions in the same packed group the unique degree-$<q$ extension of their Boolean product is
\[
 M_{ij}(Z)=
 \begin{cases}
 L_i(Z),&i=j,\\
 L_i(Z)L_j(Z)+c_ic_j\eta(Z),&i\ne j.
 \end{cases}
\]
For different groups the tensor moment is the product of one-bit moments.
\end{lemma}
\begin{proof}
$L_i$ is the unique degree-$<q$ polynomial agreeing with its bit on $K$. The product $L_iL_j$ has one degree-$q$ monomial, from multiplying both leading terms, with coefficient $c_ic_j$. All other exponents are at most $2^{g-1}+2^{g-2}<q$. Replacing $c_ic_jZ^q$ with $c_ic_jZ$ preserves all values on $K$ because $\alpha^q=\alpha$, gives degree $<q$, and proves the off-diagonal formula by uniqueness. A Boolean bit squares to itself, giving the diagonal. Tensor groups factor. The identity holds at every point of $E$, including off-grid points.
\end{proof}
For each address-register point $z_R$ ($R=A,0,1$), define the full fiber weight
\[
 w_{R,i}=\sum_{\substack{\alpha\in K\\\mathrm{low}_k(\alpha)=i}}\delta_\alpha(z_R).
\]
Each sum has exactly $q/m=32$ terms. For data-bit index $i$, write $t(i)=\lfloor i/g\rfloor$, $s(i)=i\bmod g$ and put
\[
 \mu_i=L_{s(i)}(z_{t(i)}),\quad d_i=w_{1,i}+w_{0,i},\quad
 V_r=\sum_{j<m}w_{r,j}\mu_j\quad(r=0,1).
\]
For data pack $t$, let $B_t=\sum_{j\in t}d_jc_{s(j)}$, $\eta_t=z_t^q+z_t$, and
\[
 D_i=c_{s(i)}\eta_{t(i)}(B_{t(i)}+d_ic_{s(i)})
       +d_i(\mu_i+\mu_i^2).
\]
All arithmetic is in characteristic two.
\begin{theorem}[Exact adaptive-read packed extension and joint bounds]
For all valid $m$ and all $z\in E^p$,
\[
\boxed{\displaystyle
\widetilde f_m(z)=\sum_{i<m}w_{A,i}
\big[(1+\mu_i)V_0+\mu_iV_1+D_i\big].}
\]
A single uniform deterministic multitape machine computes this extension in $O(m\operatorname{polylog}m)$ bit time and $O(\log^2m)$ writable bits. The source, field encoding, output and all input scans are charged.
\end{theorem}
\begin{proof}
Over the tensor-Lagrange weights, $f=(1-X_A)X_{B_0}+X_AX_{B_1}$, so summing the three address groups yields
\[
 \widetilde f_m=\sum_iw_{A,i}\sum_j
 [w_{0,j}(\mu_j+M_{ij})+w_{1,j}M_{ij}].
\]
Writing $C_{ij}=M_{ij}+\mu_i\mu_j$, this equals
$\sum_iw_{A,i}[(1+\mu_i)V_0+\mu_iV_1+\sum_jd_jC_{ij}]$.
Lemma \ref{lem:pair} makes $C_{ij}=0$ across packs,
$C_{ii}=\mu_i+\mu_i^2$, and
$C_{ij}=c_{s(i)}c_{s(j)}\eta_{t(i)}$ for distinct indices in one pack. Summing those terms with the one-pack statistic $B_{t(i)}$ gives exactly $D_i$.
For resources, enumerate at most $q=32m$ irreducible-degree-$g$ polynomial candidates with polynomial-$g$ tests, select a trace-one quadratic extension, and compute the trace-dual basis with $O(g^2)$ stored bits. Any $w_{R,i}$ costs only 32 Lagrange indicators, each calculated in $O(g)$ field multiplications; every $\mu_i$ uses $g$ Frobenius squarings. Scan data packs once to obtain $V_0,V_1$, rewind, then scan again computing $B_t,D_i$ and the accumulator. Neither address weights nor packed data are stored as a full table. A field multiplication costs $\operatorname{poly}(g)$ bit steps with $O(g^2)$ scratch. The basis and constant many field registers use $O(g^2)=O(\log^2m)$ bits. Total bit time is $O(m\operatorname{polylog}m)$, including source reads and result serialization.
\end{proof}
The diagonal term and the rank-one off-diagonal term are essential. Naively factoring two bit reads as if independent is \textbf{FALSIFIED}: for off-grid $z$, $z^q+z$ need not vanish and $\mu_i^2$ need not equal $\mu_i$. The same incorrect formula passes grid-only tests, where both corrections vanish.

\section{TESTS AND OUTCOME}
The full reproducibility bundle contains a standard-library finite-field implementation, seed $202610150840+$batch, original frozen proposals, eight distinct JSON checkpoints, code SHA-256
\texttt{19d9ebf2eea3cbebe569e0f037944bc672cca1a8a61cfea727480c6a21d514d6}, and an independent tensor-Lagrange reference.
The eight batches made respectively 2594, 7520, 37, 4, 25, 20, 60 and 24 assertions (10,284 total) with zero mismatches. They cover field and dual-basis tests; exact same-pack pair moments at off-grid points; $m=2,4,8,16,32$ with $q=256$ and the strong degree condition; four complete $16^4$ packed-grid reference sums (interpolation identity only); explicit off-grid failures of naive independence; black-box all-nonzero coefficients; on-grid literal indexed-read execution; and ignored trailing-source controls. Final calls were separately time-bounded to 30 seconds and all raw result files contain the code SHA and seed. Finite tests do not certify the asymptotic theorem.

\section{EXACT NONRELATIVIZATION BOUNDARY AND NEXT TARGET}
For an arbitrary Boolean oracle $Q:K^p\to\{0,1\}$ with no source representation, at an off-grid query $z$ with every $z_i\notin K$, every coefficient $\prod_i\delta_{\alpha_i}(z_i)$ is nonzero. Fewer than $q^p=2^a$ Boolean oracle-table queries cannot determine its canonical extension, since two tables differing at an unqueried location give different values. This is an exact black-box query lower bound, \emph{not} an ordinary computation lower bound. The source-specific two-read identity itself relativizes on its declared source class; the needed arbitrary-decider-to-source and complete-graph transformations remain \textbf{OPEN}. Providing an algebraic oracle value at an off-grid point defeats this query lower bound, so nonalgebrization remains independently \textbf{OPEN}. Williams's full space theorem is nonrelativizing while its Theorem 5.1 query-length-restricted oracle variant survives; both have different scope from this source theorem.

For $m=N^{3/5}$, a direct schedule of $M$ calls using this evaluator pays $M N^{3/5}\operatorname{polylog}N$ time. An $N^{11/10}$ allocation would only finance $M\le N^{1/2}/\operatorname{polylog}N$ under this method. Neither Williams nor this theorem supplies a bound on $M$ or the additional full-graph overhead, and the known space simulation does not preserve near-original time or provide arbitrary random-access coverage. LSPV-1's general all-input contraction and LSPV-3's universal compatibility obstruction are unaffected.

COMPLETIONCUT1 and LIVEPREFIX1 distinguish an unpredictable next block before disclosure from an offline SAT instance after disclosure. Freshness, waiting, hashes and source access do not require an offline decider to pre-guess. No cryptographic hardness is invoked.
The root \texttt{main.tex} remains truncated with a literal placeholder and lacks balanced LSPV-4 ownership markers; it was not edited. Existing schedules were not changed.

\noindent\textbf{One next target:} extend the moment identity to a third data-dependent indexed read with explicitly charged repeated-index coincidence partitions. Either prove a fixed-depth, near-linear/polylog-space closure theorem or preserve the smallest honest source counterexample.

\begin{thebibliography}{9}
\bibitem{Wil} R. R. Williams, \emph{Simulating Time With Square-Root Space}, STOC 2025, \url{https://people.csail.mit.edu/rrw/time-vs-space.pdf}.
\bibitem{CM} J. Cook and I. Mertz, \emph{Tree Evaluation Is in Space $O(\log n\log\log n)$}, SIAM J. Comput., 2025, DOI:10.1137/24M1690126.
\bibitem{HPR} A. Henzinger, E. Pyne, S. Ragavan, \emph{Catalytic Tree Evaluation From Matching Vectors}, 2026, arXiv:2602.14320.
\bibitem{LN} R. Lidl, H. Niederreiter, \emph{Finite Fields}, Cambridge University Press.
\end{thebibliography}
\end{document}
% END_OF_VERIFIED_REPORT::LSPV-4::ADAPTIVEMOMENT1::2026-10-15_0840KST